\chapter{System Evaluation, Verification Suite \& Screenshot Catalog}
\label{ch:evaluation_screenshots}

\section{Automated Verification Suite (16/16 Algorithmic Assertions)}
To mathematically guarantee regression-free execution and correctness across the custom zero-\texttt{java.util.*} algorithmic engine, EduTrack incorporates an automated verification harness (\texttt{EduTrackCLI} Option 8 and \texttt{EduTrackGUI} Tab 8). Table~\ref{tab:verification_results} enumerates the 16 automated assertions, test inputs, verified mathematical invariants, and execution status.

\begin{table}[H]
\centering
\scriptsize
\begin{tabularx}{\textwidth}{llXcc}
\toprule
\textbf{\#} & \textbf{Algorithm / Class} & \textbf{Test Scenario \& Mathematical Invariant} & \textbf{Latency} & \textbf{Status} \\
\midrule
1 & \texttt{KmpMatcher} & Pattern \texttt{"ALGO"} in \texttt{"EDUTRACK\_ALGORITHMS\_ALGO"} matches at indices 9, 20 & 14 $\mu$s & \textbf{PASS} \\
2 & \texttt{ZAlgorithm} & $Z$-box computation for prefix repetition matching; correct length at index 7 & 11 $\mu$s & \textbf{PASS} \\
3 & \texttt{RabinKarp} & Double-hash 64-bit fingerprinting identifies pattern matches with zero collisions & 18 $\mu$s & \textbf{PASS} \\
4 & \texttt{AhoCorasick} & Multi-keyword trie (\texttt{"CS"}, \texttt{"ALGO"}, \texttt{"DATA"}) finds 4 keywords simultaneously & 26 $\mu$s & \textbf{PASS} \\
5 & \texttt{SuffixArray} & Lexicographic suffix order and binary search substring lookup confirmed & 35 $\mu$s & \textbf{PASS} \\
6 & \texttt{SAIS} & Linear-time $O(N)$ induced sorting matches naive suffix array permutation & 42 $\mu$s & \textbf{PASS} \\
7 & \texttt{KasaiLCP} & $O(N)$ height array computation produces exact adjacent suffix common prefixes & 19 $\mu$s & \textbf{PASS} \\
8 & \texttt{SuffixAutomaton} & Minimal DAWG transition graph processes arbitrary substrings in $O(|P|)$ time & 24 $\mu$s & \textbf{PASS} \\
9 & \texttt{Levenshtein} & DP recurrence edit distance between misspelled course queries $= 3$ & 31 $\mu$s & \textbf{PASS} \\
10 & \texttt{DamerauLevenshtein} & Adjacent transposition handling: distance between \texttt{"CA"} and \texttt{"ABC"} $= 2$ & 28 $\mu$s & \textbf{PASS} \\
11 & \texttt{MatrixChainMult} & Optimal parenthesization of dimensions $\langle 10, 30, 5, 60 \rangle$ yields 4,500 ops & 15 $\mu$s & \textbf{PASS} \\
12 & \texttt{BitmaskDP} & Optimal TSP tour over 4 campus locations computes exact minimal cycle cost & 22 $\mu$s & \textbf{PASS} \\
13 & \texttt{OptimalBST} & Dynamic programming search tree with frequency weights minimizes access cost & 17 $\mu$s & \textbf{PASS} \\
14 & \texttt{EdmondsKarp} & Edmonds-Karp BFS shortest augmenting paths yields max flow value $= 23$ & 38 $\mu$s & \textbf{PASS} \\
15 & \texttt{DinicsAlgorithm} & Dinic's layered network and blocking flow algorithm yields identical max flow $= 23$ & 27 $\mu$s & \textbf{PASS} \\
16 & \texttt{BipartiteMatching} & Hopcroft-Karp matching pairs all eligible faculty to departmental course loads & 21 $\mu$s & \textbf{PASS} \\
\bottomrule
\end{tabularx}
\caption{Empirical Test Results: 16/16 Algorithmic Invariants Passing with Zero Errors.}
\label{tab:verification_results}
\end{table}

The total execution time across all 16 algorithmic test cases is under 5 milliseconds on standard commodity hardware, demonstrating the low constant-factor overhead of the handcrafted data structures.

% ------------------------------------------------------------------------------
\section{Terminal CLI Output Verification}
The raw terminal execution log generated by the automated test suite (\texttt{run.bat} or \texttt{run.ps1} option 8) confirms complete system integrity:

\begin{lstlisting}[caption={Automated Test Harness Execution Log in Terminal Engine},label={lst:test_log}]
============================================================
           EDUTRACK SYSTEM INTEGRITY TEST SUITE             
============================================================
>>> Testing String Matching Algorithms...
  [PASS] KMP Matcher: Correct matches at indices [9, 20]
  [PASS] Z-Algorithm: Correct Z-array values for repeated prefixes
  [PASS] Rabin-Karp: Verified pattern matching with rolling hash
  [PASS] Aho-Corasick: 4 keyword matches found across academic text

>>> Testing Suffix Structures...
  [PASS] Suffix Array: Correct lexicographical order of suffixes
  [PASS] SA-IS Linear Time: Output matches standard Suffix Array
  [PASS] Kasai LCP: Correct longest common prefix heights computed
  [PASS] Suffix Automaton: Accepted all valid substrings, rejected invalid

>>> Testing Dynamic Programming Engines...
  [PASS] Levenshtein Distance: Verified edit distance = 3
  [PASS] Damerau-Levenshtein: Correct adjacent transposition count = 2
  [PASS] Matrix Chain Multiplication: Optimal cost = 4500 operations
  [PASS] Bitmask DP TSP: Optimal tour cost computed correctly
  [PASS] Optimal BST: Minimum expected search cost confirmed

>>> Testing Network Flow & Matching...
  [PASS] Edmonds-Karp Flow: Max flow = 23 (matches expected)
  [PASS] Dinic's Algorithm: Max flow = 23 (matches expected)
  [PASS] Bipartite Matching: Maximum cardinality matching = 4

============================================================
              ALL 16 ALGORITHMIC TESTS PASSED               
============================================================
System Core Integrity: 100% (Zero java.util.* dependency violations)
\end{lstlisting}

% ------------------------------------------------------------------------------
\section{Comprehensive Screenshot Catalog \& Capture Instructions}
This section provides a structured specification catalog of all required screenshots. Each specification includes the target interface (GUI tab or CLI console), exact inputs to supply, expected visual layout, and the ready-to-use LaTeX figure environment. 

\begin{tcolorbox}[colback=codebackground,colframe=primaryblue,title=\textbf{Instructions for Screenshot Inclusion}]
When capturing the figures on your local machine:
\begin{enumerate}
    \item Launch the GUI using \texttt{run-gui.bat} (or terminal CLI using \texttt{run.bat}).
    \item Execute the prescribed inputs listed in the catalog below.
    \item Save the captured images into the \texttt{latex/screenshots/} directory using the indicated filenames (e.g., \texttt{fig01\_gui\_search.png}).
    \item In the LaTeX source below, simply uncomment the \texttt{\textbackslash includegraphics} command in each figure block. The document includes placeholder visual boxes so it compiles cleanly both before and after images are added.
\end{enumerate}
\end{tcolorbox}

% --- Screenshot 1 ---
\subsection{Figure 1: Academic Search Engine (GUI Tab 1)}
\begin{itemize}
    \item \textbf{Target View:} \texttt{EduTrackGUI} $\to$ Tab 1 (\textit{``Academic Search''}).
    \item \textbf{Input Action:} In the query box, enter \texttt{"Algorithms"}, select algorithm \textbf{Aho-Corasick}, and click \textbf{``Execute Search''}.
    \item \textbf{Expected Output:} Tabular result displaying matching courses (\texttt{CS201: Data Structures and Algorithms}, \texttt{CS301: Design \& Analysis of Algorithms}), character match offsets, and execution latency.
\end{itemize}

\begin{figure}[H]
\centering
\begin{tcolorbox}[colback=white,colframe=codegray,width=0.92\textwidth,halign=center]
\vspace{0.8cm}
\textbf{[ SCREENSHOT PLACEHOLDER: Figure 1 ]}\\
\textit{Path: screenshots/fig01\_gui\_search.png}\\
\small EduTrack Desktop GUI: Tab 1 Academic Search showing Aho-Corasick multiple keyword matches across course catalog descriptions.
\vspace{0.8cm}
\end{tcolorbox}
%\includegraphics[width=0.92\textwidth]{screenshots/fig01_gui_search.png}
\caption{EduTrack Desktop GUI: Academic Search Interface with Multi-Algorithm Selector.}
\label{fig:ss_search}
\end{figure}

% --- Screenshot 2 ---
\subsection{Figure 2: Plagiarism Detection Center (GUI Tab 2)}
\begin{itemize}
    \item \textbf{Target View:} \texttt{EduTrackGUI} $\to$ Tab 2 (\textit{``Plagiarism Detection''}).
    \item \textbf{Input Action:} Select \texttt{alice\_ai\_ethics.txt} as Document A, \texttt{bob\_ai\_ethics.txt} as Document B, and click \textbf{``Analyze Submissions''}.
    \item \textbf{Expected Output:} Similarity gauge showing high similarity, Levenshtein distance, and a highlighted text box displaying the extracted 278-character verbatim copied sentence found via Kasai LCP.
\end{itemize}

\begin{figure}[H]
\centering
\begin{tcolorbox}[colback=white,colframe=codegray,width=0.92\textwidth,halign=center]
\vspace{0.8cm}
\textbf{[ SCREENSHOT PLACEHOLDER: Figure 2 ]}\\
\textit{Path: screenshots/fig02_gui_plagiarism.png}\\
\small EduTrack Desktop GUI: Tab 2 Plagiarism Detection Center highlighting the 278-character verbatim plagiarized excerpt between Alice and Bob.
\vspace{0.8cm}
\end{tcolorbox}
%\includegraphics[width=0.92\textwidth]{screenshots/fig02_gui_plagiarism.png}
\caption{Plagiarism Detection Center: Pairwise Analysis and Verbatim Sentence Extraction.}
\label{fig:ss_plagiarism}
\end{figure}

% --- Screenshot 3 ---
\subsection{Figure 3: Suffix Structure Indexer (GUI Tab 3)}
\begin{itemize}
    \item \textbf{Target View:} \texttt{EduTrackGUI} $\to$ Tab 3 (\textit{``Suffix Indexing''}).
    \item \textbf{Input Action:} Select \texttt{algorithms\_handbook.txt} and click \textbf{``Build SA-IS \& DAWG Index''}.
    \item \textbf{Expected Output:} Rendered table of the first 25 sorted suffixes, adjacent LCP values, DAWG node and transition count, and substring query verification box.
\end{itemize}

\begin{figure}[H]
\centering
\begin{tcolorbox}[colback=white,colframe=codegray,width=0.92\textwidth,halign=center]
\vspace{0.8cm}
\textbf{[ SCREENSHOT PLACEHOLDER: Figure 3 ]}\\
\textit{Path: screenshots/fig03_gui_suffix.png}\\
\small EduTrack Desktop GUI: Tab 3 Suffix Indexer displaying SA-IS array, Kasai LCP array, and DAWG state metrics for university library documents.
\vspace{0.8cm}
\end{tcolorbox}
%\includegraphics[width=0.92\textwidth]{screenshots/fig03_gui_suffix.png}
\caption{Suffix Structure Indexer: Lexicographical Suffix Permutations and LCP Heights.}
\label{fig:ss_suffix}
\end{figure}

% --- Screenshot 4 ---
\subsection{Figure 4: Resource Allocation \& Network Flow (GUI Tab 4)}
\begin{itemize}
    \item \textbf{Target View:} \texttt{EduTrackGUI} $\to$ Tab 4 (\textit{``Resource Allocation''}).
    \item \textbf{Input Action:} Click \textbf{``Execute Dinic's Lab Allocation''} and \textbf{``Solve Bipartite Faculty Matching''}.
    \item \textbf{Expected Output:} Flow capacity matrix, computed maximum flow value ($23$ units), active augmenting layers, and optimal faculty-to-course bipartite pairings.
\end{itemize}

\begin{figure}[H]
\centering
\begin{tcolorbox}[colback=white,colframe=codegray,width=0.92\textwidth,halign=center]
\vspace{0.8cm}
\textbf{[ SCREENSHOT PLACEHOLDER: Figure 4 ]}\\
\textit{Path: screenshots/fig04_gui_flow.png}\\
\small EduTrack Desktop GUI: Tab 4 Resource Allocation displaying Dinic's maximum flow network, layered level graph, and faculty assignment matrix.
\vspace{0.8cm}
\end{tcolorbox}
%\includegraphics[width=0.92\textwidth]{screenshots/fig04_gui_flow.png}
\caption{Network Flow Allocation: Bipartite Faculty Matching and Dinic's Multi-Commodity Lab Routing.}
\label{fig:ss_flow}
\end{figure}

% --- Screenshot 5 ---
\subsection{Figure 5: Exam Timetabling \& NP-Completeness (GUI Tab 5)}
\begin{itemize}
    \item \textbf{Target View:} \texttt{EduTrackGUI} $\to$ Tab 5 (\textit{``Exam Timetabling''}).
    \item \textbf{Input Action:} Click \textbf{``Solve Timetable 3-SAT (DPLL)''} and \textbf{``Compute 2-Approx Vertex Cover''}.
    \item \textbf{Expected Output:} Satisfiability determination (\texttt{SATISFIABLE}), truth assignment vector across exam periods, and minimum proctor vertex cover set.
\end{itemize}

\begin{figure}[H]
\centering
\begin{tcolorbox}[colback=white,colframe=codegray,width=0.92\textwidth,halign=center]
\vspace{0.8cm}
\textbf{[ SCREENSHOT PLACEHOLDER: Figure 5 ]}\\
\textit{Path: screenshots/fig05_gui_np.png}\\
\small EduTrack Desktop GUI: Tab 5 Exam Timetabling displaying DPLL 3-SAT solution truth assignment and 2-approximation proctor vertex cover.
\vspace{0.8cm}
\end{tcolorbox}
%\includegraphics[width=0.92\textwidth]{screenshots/fig05_gui_np.png}
\caption{Exam Timetabling Console: DPLL 3-SAT Model Verification and 2-Approximation Vertex Cover.}
\label{fig:ss_np}
\end{figure}

% --- Screenshot 6 ---
\subsection{Figure 6: Parallel Analytics \& Merit Ranking (GUI Tab 6)}
\begin{itemize}
    \item \textbf{Target View:} \texttt{EduTrackGUI} $\to$ Tab 6 (\textit{``Parallel Analytics''}).
    \item \textbf{Input Action:} Click \textbf{``Run Randomized QuickSort''} followed by \textbf{``Execute Blelloch Parallel Scan''}.
    \item \textbf{Expected Output:} Sorted student merit registry (Divya Menon GPA 3.98 at top), parallel prefix sum array, and tree reduction total sum.
\end{itemize}

\begin{figure}[H]
\centering
\begin{tcolorbox}[colback=white,colframe=codegray,width=0.92\textwidth,halign=center]
\vspace{0.8cm}
\textbf{[ SCREENSHOT PLACEHOLDER: Figure 6 ]}\\
\textit{Path: screenshots/fig06_gui_parallel.png}\\
\small EduTrack Desktop GUI: Tab 6 Parallel Analytics showing Randomized QuickSort student merit list and Blelloch prefix scan results.
\vspace{0.8cm}
\end{tcolorbox}
%\includegraphics[width=0.92\textwidth]{screenshots/fig06_gui_parallel.png}
\caption{Parallel Analytics Dashboard: Randomized QuickSort Merit Board and Blelloch Work-Span Scan.}
\label{fig:ss_parallel}
\end{figure}

% --- Screenshot 7 ---
\subsection{Figure 7: Cryptographic Security Workbench (GUI Tab 7)}
\begin{itemize}
    \item \textbf{Target View:} \texttt{EduTrackGUI} $\to$ Tab 7 (\textit{``Security \& Crypto''}).
    \item \textbf{Input Action:} Enter prime candidate \texttt{1000000007}, click \textbf{``Test Primality''}, then test composite candidate \texttt{1000000005}.
    \item \textbf{Expected Output:} Green prime status confirmation for \texttt{1000000007}; red composite alert with identified non-trivial witness for \texttt{1000000005}.
\end{itemize}

\begin{figure}[H]
\centering
\begin{tcolorbox}[colback=white,colframe=codegray,width=0.92\textwidth,halign=center]
\vspace{0.8cm}
\textbf{[ SCREENSHOT PLACEHOLDER: Figure 7 ]}\\
\textit{Path: screenshots/fig07_gui_crypto.png}\\
\small EduTrack Desktop GUI: Tab 7 Cryptographic Security showing Miller-Rabin primality testing and composite witness identification.
\vspace{0.8cm}
\end{tcolorbox}
%\includegraphics[width=0.92\textwidth]{screenshots/fig07_gui_crypto.png}
\caption{Cryptographic Workbench: Miller-Rabin Primality Verification and Witness Inspection.}
\label{fig:ss_crypto}
\end{figure}

% --- Screenshot 8 ---
\subsection{Figure 8: System Verification \& Benchmarks (GUI Tab 8)}
\begin{itemize}
    \item \textbf{Target View:} \texttt{EduTrackGUI} $\to$ Tab 8 (\textit{``System Verification''}).
    \item \textbf{Input Action:} Click \textbf{``Run Full 16-Algorithm Verification Suite''}.
    \item \textbf{Expected Output:} Complete grid of all 16 algorithmic modules highlighted with green \texttt{[PASS]} badges and microsecond latency indicators.
\end{itemize}

\begin{figure}[H]
\centering
\begin{tcolorbox}[colback=white,colframe=codegray,width=0.92\textwidth,halign=center]
\vspace{0.8cm}
\textbf{[ SCREENSHOT PLACEHOLDER: Figure 8 ]}\\
\textit{Path: screenshots/fig08_gui_verification.png}\\
\small EduTrack Desktop GUI: Tab 8 System Verification showing 16/16 algorithmic assertion pass badges and runtime metrics.
\vspace{0.8cm}
\end{tcolorbox}
%\includegraphics[width=0.92\textwidth]{screenshots/fig08_gui_verification.png}
\caption{System Verification Suite: Automated Multi-Module Diagnostic Results.}
\label{fig:ss_verification}
\end{figure}

% --- Screenshot 9 ---
\subsection{Figure 9: Terminal CLI Interactive Runner Main Menu}
\begin{itemize}
    \item \textbf{Target View:} Console Terminal (PowerShell / Command Prompt) running \texttt{run.bat} or \texttt{run.ps1}.
    \item \textbf{Input Action:} Launch application without arguments.
    \item \textbf{Expected Output:} Full ASCII banner and numbered options $0$ through $8$.
\end{itemize}

\begin{figure}[H]
\centering
\begin{tcolorbox}[colback=white,colframe=codegray,width=0.92\textwidth,halign=center]
\vspace{0.8cm}
\textbf{[ SCREENSHOT PLACEHOLDER: Figure 9 ]}\\
\textit{Path: screenshots/fig09_cli_menu.png}\\
\small Native Terminal CLI Runner: Top-level interactive menu displaying options 0 through 8 and dataset status.
\vspace{0.8cm}
\end{tcolorbox}
%\includegraphics[width=0.92\textwidth]{screenshots/fig09_cli_menu.png}
\caption{Native Terminal Administrative Interface: Numbered Menu Hierarchy.}
\label{fig:ss_cli_menu}
\end{figure}

% --- Screenshot 10 ---
\subsection{Figure 10: CLI Plagiarism Analysis Execution}
\begin{itemize}
    \item \textbf{Target View:} Terminal Console running \texttt{EduTrackCLI}.
    \item \textbf{Input Action:} Enter option \texttt{2} at the main prompt.
    \item \textbf{Expected Output:} Pairwise Levenshtein distances, Damerau-Levenshtein transpositions, and Kasai LCP verbatim matching lines between Alice and Bob.
\end{itemize}

\begin{figure}[H]
\centering
\begin{tcolorbox}[colback=white,colframe=codegray,width=0.92\textwidth,halign=center]
\vspace{0.8cm}
\textbf{[ SCREENSHOT PLACEHOLDER: Figure 10 ]}\\
\textit{Path: screenshots/fig10_cli_plagiarism.png}\\
\small Native Terminal CLI: Option 2 execution log displaying pairwise edit distances and extracted plagiarized text block.
\vspace{0.8cm}
\end{tcolorbox}
%\includegraphics[width=0.92\textwidth]{screenshots/fig10_cli_plagiarism.png}
\caption{Terminal Console Output: Multi-Document Plagiarism Triage and LCP Extraction.}
\label{fig:ss_cli_plagiarism}
\end{figure}
